% Figure 10 — evidence provenance as a weighted signal field. Hollow marks are
% declarations; filled marks are work residue gathered for another purpose.
\begin{figure}[H]
\centering
\begin{tikzpicture}[pd figure,x=.89cm,y=1cm]
  \foreach \x/\lab in {3.0/identity,4.65/skill,6.30/file,7.95/note,9.60/episode,11.25/load}{
    \draw[pd guide] (\x,.55)--(\x,4.15);
    \node[pd axis label,anchor=south,rotate=35] at (\x,4.24) {\lab};
  }
  \node[pd panel title,anchor=east] at (1.70,3.15) {candidate A};
  \node[pd panel title,anchor=east] at (1.70,1.75) {candidate B};
  \draw[pd hairline] (2.05,3.15)--(12.05,3.15);
  \draw[pd hairline] (2.05,1.75)--(12.05,1.75);

  % Push evidence: equal, cheap, and visibly hollow.
  \foreach \y in {3.15,1.75}{
    \node[circle,draw=hhamber,line width=.8pt,fill=hhpaper,inner sep=2.25pt] at (3.0,\y) {};
    \node[circle,draw=hhamber,line width=.8pt,fill=hhpaper,inner sep=2.25pt] at (4.65,\y) {};
  }
  % Pull evidence: A has a file trace but no completed migration episode; B has
  % a full residue chain plus load, which also prevents blind over-routing.
  \node[pd datum] at (6.30,3.15) {};
  \node[pd caution datum] at (9.60,3.15) {};
  \foreach \x in {6.30,7.95,9.60,11.25}{\node[pd focus datum] at (\x,1.75) {};}
  \draw[pd caution rule] (9.28,2.83)--(9.92,3.47);
  \draw[pd caution rule] (9.92,2.83)--(9.28,3.47);

  \draw[pd hairline] (2.55,.72)--(5.10,.72);
  \node[pd axis label,text=hhamber,anchor=north] at (3.82,.60) {push: declared};
  \draw[pd focus rule] (5.85,.72)--(11.70,.72);
  \node[pd axis label,text=hhteal,anchor=north] at (8.78,.60) {pull: independent residue};

  % Evidence-weight summary: the bars are conclusions of the same signal field,
  % not a second table.
  \draw[pd guide] (12.72,.90)--(12.72,3.65);
  \draw[pd rule,line width=3.5pt,line cap=round] (12.72,3.15)--(13.62,3.15);
  \draw[pd focus rule,line width=3.5pt,line cap=round] (12.72,1.75)--(15.05,1.75);
  \node[pd axis label,anchor=west] at (13.82,3.15) {small verified weight};
  \node[pd direct label,text=hhteal,anchor=west,font=\footnotesize\bfseries]
    at (13.18,1.33) {B ranks first};
\end{tikzpicture}
\caption{Push and pull evidence for the same capability claim. Hollow marks are
cheap self-declarations; filled marks are residue produced while doing work for
another purpose. A and B make the same identity and skill claims, but only B
leaves the complete file, note, episode, and current-load trail. The ranker
therefore follows independently addressable evidence rather than confidence in
self-report.}
\label{fig:pushpull}
\end{figure}
